\section{Discussion and conclusion}\label{sec:discussion}
For the specified inspection task the outputs are a spatial cover, individual physically admissible examples, and a decision about another measurement. The two theorems use the same material law in different ways. Across positions the invariant object is the fixed propagation operator and the changing quantities are a few local constitutive traces. Across illuminations the invariant object is the unknown common material and the changing quantity is the applied field. Receiver-null currents account for both the failure of a positive-state test at certain positions and the limited information available from a single direction.

Three distinctions affect how these results can be used. First, a smaller position cover can be purchased with a tighter drawing or with a new measurement. Reporting the two effects separately reveals how much the measurement actually contributes. In particular, the dramatic full-field reductions in the original seven-wave experiment coexist with modest improvements after an eighth wave. Second, illumination power must be specified on the admissible source dictionary. One port used for longer and two phase-controlled ports driven together have different capabilities even when their total incident fluence is equal. The exact two- and four-port examples quantify this difference for the defined apertures and error norm. Third, a conditional quadratic response set is an outer description. A poor lower distance can express an inadequate physical relaxation, an excessively weak illumination, or an actual shared reading. Physically compatible pairs distinguish the last mechanism from the first two.

The known part outline, its allowable material values and the background bound all enter the conclusion. An arbitrary change of shape, loss, anisotropy, an out-of-plane field or an uncertain receiver gain changes the question. The three-state construction would still define a useful output for some of these problems, but the constitutive form and error bounds would have to be justified for the new model. In an experiment with random errors that diminish upon averaging, the time-sharing obstruction in \eqref{eq:transcript} is not an exposure-complexity bound. The numerical advantage in applying a rank-four update requires many changes sharing one prepared state. It should not be extrapolated to unrelated arrays or to a scan in which every pixel changes independently.

The most useful immediate output for a fixed array is an answer to a concrete physical question. A location outside the cover is inconsistent with the full nonlinear readings for the stated material tolerance. Inside the cover a realized material can be reported explicitly. When two specifications remain viable, the smallest useful energy or simultaneous port count can sometimes be bracketed before transmission. The laboratory results identify a further experimental requirement: a loaded-system error bound is needed before the same interpretation can be given to an actual inspection.

\subsection*{Data availability}
The numerical material sets, source commands saved before acquisition, noisy fields, one-sided enclosures and program listings accompany this work. The Institut Fresnel and University of Manitoba observations were obtained from the original measurement collections \cite{fresnel,manitoba}. The computations on these measurements use the common port mask specified in Section~\ref{sec:lab}.
\subsection*{Acknowledgements}
OpenAI Codex assisted with mathematical and numerical analysis, programming and manuscript preparation. Google Gemini contributed independent technical review. Responsibility for the scientific claims and the verification of the derivations and numerical bounds rests with the authors.
